\src{8.251 Test 1 Q2, (a)--(b), 20 分; 2007 年 3 月 21 日原卷第 1、3 页.}
\begin{exercise}
复平面有识别 $z\sim\ee^{2\pi\ii/3}z$. 原卷附图的 $A,B$ 均在正实轴上, $0<r_A<r_B$, 该轴作为基本域的一边, 基本域应位于其上方.
(a) 绘出第二边并标明它与坐标轴的夹角.
(b) 有多少条测地线连接 $A,B$? 绘出这些测地线, 包括它们在基本域内的表示. 下图保留原题点的先后位置, 未给定的距离不擅自补成数值条件.
\begin{center}
\begin{tikzpicture}[scale=.85,>=Stealth]
 \draw[->](-.5,0)--(4.1,0)node[right]{实轴};
 \draw[->](0,-.4)--(0,2.4)node[above]{虚轴};
 \fill(2,0)circle(2pt)node[below]{$A$};
 \fill(3,0)circle(2pt)node[below]{$B$};
 \node at(2.1,1.7){待补基本域};
\end{tikzpicture}
\end{center}
\end{exercise}
\begin{solution}
生成元旋转 $120^\circ$, 基本域为 $0\leq\arg z<120^\circ$, 第二边与正实轴夹角 $120^\circ$, 两边按旋转识别. 射线上 $A$ 的副本为 $A'=r_A\ee^{2\pi\ii/3}$, $B$ 的副本为 $B'=r_B\ee^{2\pi\ii/3}$.

固定 $A$ 的覆盖空间原像, $B$ 有三个群像 $B,\ee^{2\pi\ii/3}B,\ee^{4\pi\ii/3}B$. 欧氏覆盖空间中到每个像有唯一一条直线段, 均不经过原点, 降到商空间后为三条不同的局部测地线. 基本域内可画成 $AB,AB',A'B$, 其中最后一条来自下半平面段的旋转代表. 只数其中最短的 $AB$ 会把“全部测地线”误解为“全局最短路径”.

\begin{center}
\begin{tikzpicture}[scale=.9,>=Stealth]
 \fill[structurecolor!7](0,0)--(3.5,0)arc(0:120:3.5)--cycle;
 \draw[->](-.2,0)--(3.9,0);
 \draw[thick](0,0)--(120:3.5);
 \draw(1,0)arc(0:120:1);\node at(.4,1.2){$120^\circ$};
 \coordinate(A)at(2,0);\coordinate(B)at(3,0);
 \coordinate(AA)at(120:2);\coordinate(BB)at(120:3);
 \draw[very thick,structurecolor](A)--(B);
 \draw[very thick,structurecolor](A)--(BB);
 \draw[very thick,structurecolor,dashed](AA)--(B);
 \fill(A)circle(2pt)node[below]{$A$};\fill(B)circle(2pt)node[below]{$B$};
 \fill(AA)circle(2pt)node[left]{$A'$};\fill(BB)circle(2pt)node[left]{$B'$};
\end{tikzpicture}
\end{center}
为进一步核对, 三条长度分别是 $r_B-r_A$ 和两份
$\sqrt{r_A^2+r_B^2+r_Ar_B}$. 它们的端点在商空间相同, 路径却不同; 经过锥尖点的折线不是本题这三条光滑提升的测地线, 且也不是更短的连接. 绘图用任意示意半径 2、3, 只表达 $r_A<r_B$, 不把这些画图值写成原题测量数据.
\end{solution}
